\section{\personFloer{}链映射的复合法则}
\label{sec:composition-law-floer-chain-maps}

本节考察\personFloer{}同态
\begin{equation}
 h_{\alpha\beta}:HF_*(H_\alpha)\longrightarrow HF_*(H_\beta)
\label{eq:19-floer-homology-continuation-map}
\end{equation}
的复合法则
\[
 h_{\alpha\gamma}=h_{\beta\gamma}\circ h_{\alpha\beta}.
\]
现在重新考察方程 (18.5.42)。研究\personFloer{}模空间的一个关键分析事实是存在先验能量
上界；下一节将对此作说明。

考虑哈密顿量三元组
\[
 (H_\alpha,H_\beta,H_\gamma)
\]
以及同伦 $\mathcal H_1=\{H_1(s)\}_{0\leq s\leq1}$、
$\mathcal H_2=\{H_2(s)\}_{0\leq s\leq1}$，满足
\[
 H_1(0)=H_\alpha,\quad H_1(1)=H_\beta=H_2(0),\quad H_2(1)=H_\gamma.
\]
定义它们的拼接
$\mathcal H_1\#\mathcal H_2=\{H_3(s)\}_{0\leq s\leq1}$ 为
\[
 H_3(s)=
 \begin{cases}
  H_1(2s),&0\leq s\leq\frac12,\\
  H_2(2s-1),&\frac12\leq s\leq1.
 \end{cases}
\]

注意，由于截断函数 $\rho$ 的选取，延拓方程 (18.5.42) 在区域 $|\tau|>R$ 上是
\emph{自治}的，即在 $\tau$ 的平移下不变。给定三元组
$(H_\alpha,H_\beta,H_\gamma)$ 时，这一事实使我们能够分别粘合对应于
$(H_\alpha,H_\beta)$ 与 $(H_\beta,H_\gamma)$ 的两个这种方程的解。

下面作更精确的说明。给定一对拉长函数
\[
 \rho=(\rho_1,\rho_2)
\]
和正数 $R>0$，定义 $\mathcal H_1\#\mathcal H_2$ 的拉长同伦
\[
 \mathcal H_1\#_{(\rho;R)}\mathcal H_2
 =\{H_{(\rho;R)}(\tau)\}_{-\infty<\tau<\infty}
\]
如下：
\[
 H_{(\rho;R)}(\tau,t,x)=
 \begin{cases}
 H_1(\rho_1(\tau+2R),t,x),&\tau\leq0,\\
 H_2(\rho_2(\tau-2R),t,x),&\tau\geq0.
 \end{cases}
\]
注意，对分别依赖于拉长函数 $\rho_1,\rho_2$ 与同伦
$\mathcal H_1,\mathcal H_2$ 的某些充分大的 $R_1,R_2>0$，有
\[
 H_{(\rho;R)}\equiv
 \begin{cases}
 H_\alpha,&\tau\leq-(R_1+2R),\\
 H_\beta,&-R\leq\tau\leq R,\\
 H_\gamma,&\tau\geq R_2+2R.
 \end{cases}
\]
\emph{特别地，即使通常的粘合同伦 $\mathcal H_1\#\mathcal H_2$ 未必光滑，这个拉长同伦
也总是光滑的}。类似地定义 $\mathfrak j_1\#_{(\rho;R)}\mathfrak j_2$。

对拉长同伦
$(\mathfrak j_1\#_{(\rho;R)}\mathfrak j_2,
\mathcal H_1\#_{(\rho;R)}\mathcal H_2)$，考虑相应的扰动\personCauchy{}--\personRiemann{}方程
\[
 \left\{
 \begin{aligned}
 &\frac{\partial u}{\partial\tau}
 +J_3^{\rho(\tau)}
 \left(\frac{\partial u}{\partial t}-X_{H_3^{\rho(\tau)}}(u)\right)=0,\\
 &\lim_{\tau\to-\infty}u(\tau)=z^-,\qquad
 \lim_{\tau\to\infty}u(\tau)=z^+,
 \end{aligned}
 \right.
\]
并附加条件~\eqref{eq:19-chain-homotopy-class-condition}。

设 $u_1,u_2$ 分别是与 $\rho_1,\rho_2$ 对应的方程 (18.5.44) 的解。在区域
$-R\leq\tau\leq R$ 上用一个单位分解作适当的固定插值，并在两端令预粘合映射为
\[
 u_1\#_Ru_2(\tau,t)=
 \begin{cases}
 u_1(\tau+2R,t),&\tau\leq-R,\\
 u_2(\tau-2R,t),&\tau\geq R.
 \end{cases}
\]
若 $R_0$ 充分大，则赋值
\[
 (u_1,u_2,R)\longmapsto u_1\#_Ru_2
\]
定义从
\[
\begin{aligned}
 &\mathcal M(\mathfrak j_1,\mathcal H_1;[z_1,w_1],[z_2,w_2])
 \times\mathcal M(\mathfrak j_2,\mathcal H_2;[z_2,w_2],[z_3,w_3])\\
 &\hspace{8em}\times(R_0,\infty)
\end{aligned}
\]
到其像的微分同胚。

记 $\overline\partial_{(\mathcal H,\mathfrak j;\rho)}$ 为相应的扰动\personCauchy{}--\personRiemann{}算子
\[
 u\longmapsto\frac{\partial u}{\partial\tau}
 +J_3^{\rho(\tau)}
 \left(\frac{\partial u}{\partial t}-X_{H_3^{\rho(\tau)}}(u)\right),
\]
它作用在满足渐近条件 $u(\pm\infty)=z^\pm$ 和固定同伦条件
$[u]=C\in\pi_2(z^-,z^+)$ 的映射上。在适当的 $u$ 的 $W^{1,p}$ 空间（$p>2$）中，
把 $u_1\#_Ru_2$ 扰动一个小于它成为真解所差误差的量，也就是小于加权 $L^p$-范数
\[
 \bigl\|\overline\partial_{(\mathcal H,\mathfrak j;\rho)}
 (u_1\#_{(\rho;R)}u_2)\bigr\|_p,
\]
便可在 $u_1\#_Ru_2$ 附近构造唯一的真解。利用第 15.5 节的粘合定理，略用记号，
仍把这个真解记作 $u_1\#_Ru_2$。相应映射于是给出嵌入
\[
\begin{aligned}
 &\mathcal M(\mathfrak j_1,\mathcal H_1;[z_1,w_1],[z_2,w_2])
 \times\mathcal M(\mathfrak j_2,\mathcal H_2;[z_2,w_2],[z_3,w_3])
 \times(R_0,\infty)\\
 &\quad\longrightarrow
 \mathcal M(\mathfrak j_1\#_{(\rho;R)}\mathfrak j_2,
 \mathcal H_1\#_{(\rho;R)}\mathcal H_2;[z_1,w_1],[z_3,w_3]).
\end{aligned}
\]
特别地，当
\[
 \mu_{H_\beta}([z_2,w_2])-\mu_{H_\alpha}([z_1,w_1])
 =\mu_{H_\gamma}([z_3,w_3])-\mu_{H_\beta}([z_2,w_2])=0
\]
时，$\mathcal M(\mathfrak j_1,\mathcal H_1;[z_1,w_1],[z_2,w_2])$ 与
$\mathcal M(\mathfrak j_2,\mathcal H_2;[z_2,w_2],[z_3,w_3])$ 都紧致，因而由有限多个点
组成。此外，上述嵌入的像穷尽模空间
\[
 \mathcal M(\mathfrak j_1\#_{(\rho;R)}\mathfrak j_2,
 \mathcal H_1\#_{(\rho;R)}\mathcal H_2;[z_1,w_1],[z_3,w_3])
\]
的“端”，其紧化的边界由断裂流线
\[
 u_1\#_{(\rho;\infty)}u_2=u_1\#_\infty u_2
\]
组成。这就证明了如下粘合恒等式。

\begin{proposition}
\label{prop:composition-gluing-chain-maps}
存在 $R_0>0$
\footnote{译者注：原书此式右端把两个链映射的复合次序写成
$h_{(\mathcal H_1,\mathfrak j_1;\rho_1)}\circ
h_{(\mathcal H_2,\mathfrak j_2;\rho_2)}$，但这与陈述中的定义域、陪域不匹配，
也与本节开头的 $h_{\alpha\gamma}=h_{\beta\gamma}\circ h_{\alpha\beta}$ 相反；
译文按后者改正。}
，使得对任意 $R\geq R_0$，作为从 $CF_*(H_\alpha)$ 到
$CF_*(H_\gamma)$ 的链映射，有
\[
 h_{(\mathcal H_1,\mathfrak j_1)\#_{(\rho;R)}
 (\mathcal H_2,\mathfrak j_2)}
 =h_{(\mathcal H_2,\mathfrak j_2;\rho_2)}
 \circ h_{(\mathcal H_1,\mathfrak j_1;\rho_1)}.
\]
\end{proposition}

这里提醒读者，同伦 $\mathcal H_1\#_{(\rho;R)}\mathcal H_2$ 本身是粘合同伦
$\mathcal H_1\#\mathcal H_2$ 的一个拉长。此命题于是给出同调上的复合法则
$h_{\alpha\gamma}=h_{\beta\gamma}\circ h_{\alpha\beta}$。
